\documentclass[10pt]{article}

\usepackage[margin=0.75in]{geometry}
\usepackage{amsmath,amsthm,amssymb,bm}
\usepackage{xcolor}
\usepackage{cancel}
\usepackage{graphicx}
\usepackage{changepage}
\usepackage{circuitikz}
\usepackage{pgfplots}
\usepackage{minted}
\usepackage{physics}
\usepackage{tikz}
\usepackage{tikz-3dplot}
\usepackage{hyperref}
\usepackage{siunitx}
\usepackage{fontspec}
\usepackage{relsize}
\usepackage{subfig}
\usepackage{todonotes}
\usepackage{contour}
\usepackage{multicol, multirow, booktabs}
\usepackage[breakable]{tcolorbox}
\usepackage[inline]{enumitem}

\theoremstyle{definition}
\newtheorem{problem}{Problem}
\newtheorem{soln}{Solution}

\pgfplotsset{compat=newest}
\usetikzlibrary{lindenmayersystems}
\usetikzlibrary{arrows}
\usetikzlibrary{calc}
\usetikzlibrary{positioning, fit}
\usetikzlibrary{3d, perspective}
\usetikzlibrary{math}
\usetikzlibrary{fadings}
\usetikzlibrary{angles,quotes} 
\usetikzlibrary{decorations.markings}

\definecolor{incolor}{HTML}{303F9F}
\definecolor{outcolor}{HTML}{D84315}
\definecolor{cellborder}{HTML}{CFCFCF}
\definecolor{cellbackground}{HTML}{F7F7F7}
\newcommand{\ui}{\hat{i}}
\newcommand{\uj}{\hat{j}}
\newcommand{\uk}{\hat{k}}
\newcommand{\ux}{\hat{\mathbf{x}}}
\newcommand{\uy}{\hat{\mathbf{y}}}
\newcommand{\uz}{\hat{\mathbf{z}}}
\newcommand{\uv}[1]{\hat{\bv{{#1}}}}
\newcommand{\pr}[1]{{#1^\prime}}
\newcommand{\pd}[2][]{{\frac{\partial^{#1}}{\partial {{#2}^{#1}}}}}
\newcommand{\dif}[2][]{{\frac{d^{#1}}{d{{#2}^{#1}}}}}
\newcommand{\grd}{\nabla}

\pgfdeclarelayer{background}  
\pgfsetlayers{background,main}
\AtBeginDocument{\RenewCommandCopy\qty\SI}

\makeatletter
\newcommand{\boxspacing}{\kern\kvtcb@left@rule\kern\kvtcb@boxsep}
\makeatother
\newcommand{\prompt}[4]{
    \ttfamily\llap{{\color{#2}[#3]:\hspace{3pt}#4}}\vspace{-\baselineskip}
}

\newcommand{\thevenin}[2]{
  \begin{center}
    \begin{circuitikz} \draw
      (0,0) -- (2,0) to[battery1, l_=$V_{Th}\eq#1$] (2,2) 
      to[resistor, l_=$R_{Th}\eq#2$] (0,2)
      ;
      \draw [o-] (-.07,2.079);
      \draw [o-] (-.07,0.079);
    \end{circuitikz}
  \end{center}
}

\newcommand{\norton}[2]{
  \begin{center}
    \begin{circuitikz} \draw
      (0,0) -- (3,0) to[american current source, l_=$I_{N}\eq#1$] (3,2) -- (0,2) (2,0)
      to[resistor, l=$R_{N}\eq#2$] (2,2)
      ;
      \draw [o-] (-.07,2.079);
      \draw [o-] (-.07,0.079);
    \end{circuitikz}
  \end{center}
}

\DeclareMathOperator{\Div}{div}
\DeclareMathOperator{\Curl}{curl}
\DeclareMathOperator{\sgn}{sgn}

\newcommand{\highlight}[1]{\colorbox{yellow}{$\displaystyle #1$}}
\newcommand{\ti}[1]{\widetilde{#1}}
\renewcommand{\div}[1]{\bv{\nabla}\cdot #1}
\renewcommand{\curl}[1]{\bv{\nabla}\times #1}

\newfontface{\Kaufmann}{Kaufmann}
\DeclareTextFontCommand{\kf}{\Kaufmann}
\newcommand{\scriptr}{\fontsize{12pt}{12pt}\kf{r}\fontsize{10pt}{10pt}}

\newfontface{\KaufmannB}{Kaufmann Bold}
\DeclareTextFontCommand{\kfb}{\KaufmannB}
\newcommand{\bscriptr}{\fontsize{12pt}{12pt}\kfb{r}\fontsize{10pt}{10pt}}
\newcommand{\justif}[2]{&{#1}&\text{#2}}
\newcommand{\bv}[1]{\bm{\mathrm{#1}}}

\title{Physics 4700H: Assignment III}
\author{Jeremy Favro (0805980) \\ Trent University, Peterborough, ON, Canada}
\date{\today}

\begin{document}
\maketitle

% PROBLEM 1
\begin{problem} Consider, again, the case of $N$ tosses of a fair coin.\\
On PS2, you considered the multiplicity function describing the number of ways a value of
$X = N_H -N_T$ can be obtained. Using the more approximate form of Stirling's approximation
you showed that $\Omega(N, X)$ is approximately Gaussian when $N \gg 1$ and $\abs{X}/N\ll1$.
However, the form you found is not normalized.
If one uses the less approximate form of Stirling's approximation, one obtains
$$\Omega(N,X)\approx\sqrt{\frac{2}{\pi N}}2^Ne^{-X^2/2N}$$
Show that in this case, the sum of the probability $P (N, X) = \Omega(N, X)/2^N$ over all possible
$X$ values is one. In the limit of $N$ is large, the sum can be converted to an integral, but do
think carefully about the limits of the sum and the conversion to the integral, and articulate
clearly the approximations you are making.
\end{problem}
\begin{soln}
Just to write it out, $$P (N, X) = \Omega(N, X)/2^N\approx \sqrt{\frac{2}{\pi N}}e^{-X^2/2N}.$$
We want to show that
$$1=\sqrt{\frac{2}{\pi N}}\sum_{X=-N}^{N}e^{-X^2/2N}.$$
To convert this to an integral we need to have $N$ large enough that $dX=2$ is insignificant, i.e. that
$P(N,X)=P_c(N,X)dX$ in the continuous case. This forces that, dividing through by $dX$, 
$$\frac{1}{2}P(N,X)=P_c(N,X)=\frac{1}{\sqrt{2\pi N}}e^{-X^2/2N}.$$
Integrating this
$$\frac{1}{\sqrt{2\pi N}}\int_{-N}^{N}e^{-X^2/2N}dX$$
we let $u=X/\sqrt{2N}\implies \sqrt{2N}du=dX$ which transforms the integral to
$$\frac{\sqrt{2N}}{\sqrt{2N\pi }}\int_{-\sqrt{N/2}}^{\sqrt{N/2}}e^{-u^2}du=\frac{1}{\sqrt{\pi}}\int_{-\sqrt{N/2}}^{\sqrt{N/2}}e^{-u^2}du$$
which is half of the error function, $\erf(x)$. So 
$$\frac{1}{\sqrt{\pi}}\int_{-\sqrt{N/2}}^{\sqrt{N/2}}e^{-u^2}du=\frac{1}{2}\left[\erf(\sqrt{N/2})-\erf(-\sqrt{N/2})\right]$$
and since the error function is odd,
$$\frac{1}{\sqrt{\pi}}\int_{-\sqrt{N/2}}^{\sqrt{N/2}}e^{-u^2}du=\erf(\sqrt{N/2}).$$
Now $\erf(x)$ asymptotes to 1 \emph{very} quickly. As an example, to 512(!) bits of precision, SageMath gives 
$$\erf(12)\approx0.999999999999999999999999999999999999999999999999999999999999999864373883079409578721969384340958242733321776671189507008762046251464309514787536774406542$$
(page cutoff for comedic effect). I'd call that normalized.
\end{soln}
\end{document}